\chapter{Introduction}
\label{ch:introduction}

\section{Motivation: two failures that look unrelated}

Modern machine learning exhibits two persistent, well-documented failure classes that
are almost always studied in isolation. The first is \emph{compositional generalisation}
failure: a model trained on known components and combinations fails to recombine them
under novel distribution, performing dramatically worse on combinations that are
structurally trivial yet unseen during training \citep{lake2018generalization,
keysers2020measuring,hupkes2023taxonomy}. The second is \emph{alignment} failure: a
system optimises competently toward objectives that diverge from what its designers
intended, whether through misspecification \citep{amodei2016concrete}, learned
optimisation pressures \citep{hubinger2019risks}, or catastrophic risk profiles
\citep{hendrycks2023overview}.

Both failure classes share a striking signature. In compositional generalisation
failure, the model's internal representations have collapsed onto the surface
statistics of the training distribution; the model has, in effect, lost the structural
content that would let it compose known parts \citep{lake2018generalization,
geirhos2020shortcut}. In alignment failure, the system's internal objective has
drifted from the specification it was trained to satisfy; the system pursues a proxy
that is coherent \emph{internally} but detached from the intended semantics
\citep{amodei2016concrete,russell2019human}. In both cases, what has failed is the
relationship between the internal structure of the system and the structure of the
world it is supposed to track.

\section{The $\sigma$-trap thesis}

This monograph develops a single claim:

\begin{quote}
\emph{Compositional generalisation failure and AI alignment failure are the same
phenomenon --- a bifurcation in the agent's internal schema coherence (the
$\sigma$-trap) --- and solving one solves the other.}
\end{quote}

We call the internal degree of consistency of an agent's representation structure its
\emph{schema coherence} $\sigma$ (Chapter~4 gives the formal definition). The claim is
that under gradient-based learning, $\sigma$ behaves like an order parameter of a
dynamical system: above a critical value $\sigmacrit$ learning is compositional and
robust; below it, the system is trapped in a stable equilibrium of low coherence from
which ordinary training does not escape. The $\Sigma$-Model, a minimal mechanistic
model developed in Chapter~7, exhibits exactly this bifurcation
\citep{amsyar2026sigma}. If the claim holds, alignment is not a separate engineering
problem bolted onto the top of a model: it is the same structural problem as
compositional generalisation, and interventions that raise schema coherence
($\sigma$-coupling, curriculum design) are simultaneously safety interventions.

The thesis is deliberately structured so that this strong claim is not presupposed by
the evidence reviews. Chapters~2--3 map the literature and the evidence base using
schema-coherence vocabulary strictly as a heuristic lens, and report what is found
--- including the limited explicit intersection between the compositional
generalisation and alignment literatures --- without requiring the arc claim to hold.
The claim is argued from Chapter~4 onward, where the framework, the mechanistic
model, and the interventions can carry the weight.

\section{Why a monograph}

The claim spans two research communities that rarely read each other: the
compositional generalisation community, which works on benchmarks and mechanistic
accounts of systematicity \citep{keysers2020measuring,hupkes2023taxonomy,
lake2023compositional}, and the AI safety community, which works on specification,
robustness, and existential risk \citep{bostrom2014superintelligence,
hendrycks2023overview,ji2023survey}. A unified claim needs a unified presentation:
definitions that serve both communities, a shared notation, and a continuous argument
rather than a collection of standalone papers. The monograph form makes that possible.
Underlying publications --- the scoping review, the systematic review, the empirical
$\Sigma$-Model study, and the papers that will follow --- are listed in the front
matter and archived under \texttt{thesis/publications/}; chapters adapt them into
monograph prose rather than reproducing venue formatting.

\section{Contributions}

The monograph makes the following contributions.

\begin{enumerate}
  \item \textbf{Diagnosis.} A scoping review of the AGI safety landscape (Ch.~2) and a
    systematic review of the $\sigma$-trap evidence base (Ch.~3) establish that the
    safety literature is fragmented and that no existing framework treats internal
    schema coherence as the central failure mode.
  \item \textbf{Framework.} The $\Sigma$-Align framework (Ch.~4) formalises schema
    coherence as a measurable, intervenable quantity and gives definitions for the
    $\sigma$-trap and its bifurcation.
  \item \textbf{Evidence.} A pilot study of $\sigma$-coupling interventions (Ch.~5),
    a meta-analysis quantifying the $\sigma$-trap (Ch.~6), the $\Sigma$-Model showing
    compositional generalisation failure as a bifurcation (Ch.~7), and a study linking
    the trap to mesa-optimization (Ch.~8).
  \item \textbf{Implications.} An argument that schema-coherent training is the
    optimal path to safe, long-term AGI, with direct relevance to corrigibility, CEV,
    and indirect normativity (Ch.~9).
\end{enumerate}

\section{Reader's guide}

Chapter~2 maps the landscape of AGI safety research and identifies the gaps a
schema-coherence framework could address. Chapter~3 reviews the evidence for the
$\sigma$-trap as a measurable construct. Chapter~4 builds the $\Sigma$-Align
framework. Chapters~5--8 present the empirical programme: the pilot, the
meta-analysis, the $\Sigma$-Model, and the mesa-optimization study. Chapter~9 draws
out the implications for safe AGI, and Chapter~10 concludes.

The chapters are written to be read in order, but each chapter opens with a summary of
its position in the argument and closes with a ``Relation to Other Chapters'' section,
so that readers from either community can enter at the chapter most relevant to them.
